\documentclass[a4paper,12pt]{article}
\usepackage[utf8]{inputenc}
\usepackage[german]{babel}
\usepackage{amsmath}
\usepackage{amssymb}
\usepackage{amsthm}
\usepackage{geometry}
\usepackage{graphicx}
\usepackage[hidelinks]{hyperref}
\usepackage{listings}
\usepackage{xcolor}
\usepackage{enumitem}
\usepackage{booktabs}
\usepackage{array}
\usepackage{multirow}
\usepackage{float}
\usepackage{caption}
\usepackage{subcaption}
\usepackage{siunitx}

\geometry{margin=2.5cm}

\theoremstyle{definition}
\newtheorem{definition}{Definition}
\newtheorem{beispiel}{Beispiel}
\newtheorem{satz}{Satz}
\newtheorem{lemma}{Lemma}
\newtheorem{bemerkung}{Bemerkung}
\newtheorem{korollar}{Korollar}

\setlist[itemize,1]{label=--}

% ---------- Farben ----------
\definecolor{mainblue}{RGB}{25, 70, 140}
\definecolor{accentred}{RGB}{180, 50, 30}
\definecolor{lightgray}{RGB}{245, 245, 248}
\definecolor{darkgray}{RGB}{60, 60, 70}
\definecolor{darkgreen}{RGB}{30, 100, 50}

\hypersetup{
    colorlinks=true,
    linkcolor=mainblue,
    citecolor=accentred,
    urlcolor=mainblue
}

\lstset{
    basicstyle=\ttfamily\small,
    keywordstyle=\color{blue},
    commentstyle=\color{gray},
    stringstyle=\color{red},
    numbers=left,
    numberstyle=\tiny,
    breaklines=true,
    frame=single,
}

\title{\textbf{Methanthiol als natürliche Klimanotbremse des Ozeans:}\\
\large Formale Analyse mariner Schwefelemissionen, atmosphärischer Oxidationschemie\\
und Aerosolstrahlungsantrieb im Südpolarmeer}
\author{Stephan Epp\\
\
\small \texttt{Stephan\_Epp@web.de}}
\date{14. Mai 2026}

\begin{document}

\maketitle

\begin{abstract}
Diese Arbeit analysiert formal die Entdeckung von Wohl et al.~(2024), nach der marine Emissionen von
Methanthiol~(MeSH, \(\mathrm{CH_3SH}\)) die Aerosolkühlung über dem Südpolarmeer um 30--70\,\%
erhöhen und damit eine bislang unbekannte natürliche Rückkopplungsschleife des Klimasystems
darstellen. Wir entwickeln ein mathematisch-physikalisches Rahmenwerk, das den marinen
Schwefelzyklus, die atmosphärische Oxidationschemie von MeSH, die Aerosolbildung sowie den
resultierenden Strahlungsantrieb (RF) formal beschreibt. Auf Basis von Emissionsgleichungen,
Reaktionskinetik, Aerosol-Nukleationsmodellen und Strahlungstransferformulierungen werden zentrale
Lemmas und Sätze bewiesen. Acht Matplotlib-Visualisierungen veranschaulichen die quantitativen
Ergebnisse. Ein ausführlicher Ausblick diskutiert offene wissenschaftliche Fragen sowie
Implikationen für Klimamodellierung, Geo-Engineering und Klimapolitik.\\[0.4em]
\textbf{Schlüsselwörter:} Methanthiol (MeSH), Dimethylsulfid (DMS), marines Aerosol, Südpolarmeer,
Strahlungsantrieb, Schwefelzyklus, CLAW-Hypothese, natürlicher Klimarückkopplungsmechanismus.
\end{abstract}

\tableofcontents
\newpage

% ─────────────────────────────────────────────────────────────────────────────
\section{Einführung}
% ─────────────────────────────────────────────────────────────────────────────

\subsection{Problemstellung und wissenschaftliche Relevanz}

Das Klimasystem der Erde ist ein hochkomplexes, nichtlineares System aus Wechselwirkungen zwischen
Ozean, Atmosphäre, Biosphäre und Kryosphäre. Inmitten der wachsenden anthropogenen Treibhausgaskonzentrationen
besitzen natürliche negative Rückkopplungsmechanismen eine besondere Bedeutung: Sie können die
Erwärmung partiell dämpfen und als „natürliche Notbremsen" fungieren.

Im November 2024 veröffentlichten Wohl et al.\ in \textit{Science Advances} einen Befund, der das
Verständnis des marinen Schwefelzyklus grundlegend verändert. Die Studie zeigt:

\begin{itemize}
  \item Der dominierende Weg des marinen biogenen Schwefelstoffwechsels führt \emph{nicht} durch
        Dimethylsulfid~(DMS), wie bisher angenommen, sondern zu einem wesentlichen Teil durch
        Methanthiol~(MeSH, \(\mathrm{CH_3SH}\)).
  \item Marine MeSH-Emissionen erhöhen die Sulfataerosolsäulendichte über dem Südpolarmeer um
        \textbf{30--70\,\%}.
  \item Durch den vermehrten Aerosolgehalt wird der kühlende Strahlungsantrieb verstärkt, der
        Strahlungsbias aktueller Klimamodelle wird signifikant reduziert.
  \item MeSH-Oxidation verbraucht atmosphärische Oxidanzien~(OH) und verlängert dadurch die
        Lebensdauer von DMS, was dessen Ferntransport und klimatischen Einfluss vergrößert.
\end{itemize}

Diese Arbeit zielt darauf ab, die genannten Ergebnisse in einem formalen mathematisch-physikalischen
Rahmen darzustellen, Beweise für die zentralen Aussagen zu führen und die quantitativen Ergebnisse
zu veranschaulichen.

\subsection{Stand der Forschung: Mariner Schwefelzyklus und DMS}

Das Konzept, dass marine Organismen das Klima über Schwefelverbindungen regulieren, wurde erstmals
1987 von Charlson, Lovelock, Andreae und Warren in der sogenannten \textbf{CLAW-Hypothese} formuliert
\cite{charlson1987}. Der Mechanismus sieht vor: Phytoplankton produziert Dimethylsulfoniopropionat
(DMSP); dieses wird zu DMS umgesetzt, das in die Atmosphäre emittiert wird; dort oxidiert DMS zu
Sulfat, das als Wolkenkondensationskeim (CCN) dient und damit die Wolkenalbedo und Klimakühlung
verstärkt.

Jahrelang wurde DMS als der dominante marine Schwefelemittent behandelt~\cite{lana2011,
andreae2019}. MeSH entsteht jedoch aus demselben Substrat~(DMSP) wie DMS, über einen alternativen
biochemischen Pfad~\cite{kiene1996}, und übersteigt in vielen Meeresbereichen tatsächlich die
Flüsse durch DMS. Die hohe Reaktivität von MeSH mit OH-Radikalen hat bislang zuverlässige
Messungen in Meerwasser und Atmosphäre verhindert.

\subsection{Ziel dieser Arbeit}

Ziel dieser Arbeit ist es:
\begin{enumerate}
  \item Den marinen Schwefelzyklus unter Einbezug von MeSH formal zu modellieren.
  \item Die atmosphärische Oxidationschemie von MeSH mathematisch zu beschreiben.
  \item Den resultierenden Aerosolstrahlungsantrieb quantitativ herzuleiten.
  \item Den negativen Klimarückkopplungsmechanismus formal als Satz zu beweisen.
  \item Ergebnisse anhand von acht Matplotlib-Abbildungen zu veranschaulichen.
\end{enumerate}

% ─────────────────────────────────────────────────────────────────────────────
\section{Grundlagen}
% ─────────────────────────────────────────────────────────────────────────────

\subsection{Chemische Grunddefinitionen}

\begin{definition}[DMSP]
  Dimethylsulfoniopropionat (\(\mathrm{(CH_3)_2S^+CH_2CH_2COO^-}\), DMSP) ist eine organische
  Schwefelverbindung, die von marinem Phytoplankton synthetisiert wird und als Osmoprotektant
  fungiert. DMSP ist der primäre Vorläufer sowohl von DMS als auch von MeSH im marinen
  Schwefelzyklus.
\end{definition}

\begin{definition}[DMS und MeSH]
  Dimethylsulfid~(\(\mathrm{(CH_3)_2S}\), DMS) und Methanthiol~(\(\mathrm{CH_3SH}\), MeSH) sind
  flüchtige organische Schwefelverbindungen. Beide entstehen enzymatisch aus DMSP:
  \begin{align}
    \mathrm{DMSP} &\xrightarrow{\text{DMSP-Lyase}} \mathrm{DMS} + \mathrm{Acrylat} \label{eq:dms_formation}\\
    \mathrm{DMSP} &\xrightarrow{\text{Demethylierung}} \mathrm{MMSP} \xrightarrow{} \mathrm{MeSH} + \ldots \label{eq:mesh_formation}
  \end{align}
\end{definition}

\begin{definition}[Luftsee-Fluss]
  Der Luft-See-Fluss~$F$ eines gelösten flüchtigen Gases~$X$ wird beschrieben durch:
  \[
    F_X = k_w \cdot (C_w - C_a / H_X)
  \]
  wobei $k_w$ die wasserseitige Transfergeschwindigkeit~[\si{m\,s^{-1}}], $C_w$ die
  Meerwasserkonzentration~[\si{mol\,m^{-3}}], $C_a$ die atmosphärische Konzentration und $H_X$
  die dimensionslose Henry-Konstante ist.
\end{definition}

\begin{definition}[OH-Radikale und atmosphärische Lebensdauer]
  OH-Radikale sind die wichtigsten Oxidanzien in der Troposphäre. Die atmosphärische Lebensdauer
  $\tau_X$ einer Verbindung $X$ gegenüber OH-Oxidation ist:
  \[
    \tau_X = \frac{1}{k_{X+\mathrm{OH}} \cdot [\mathrm{OH}]}
  \]
  mit der bimolekularen Ratenkonstante $k_{X+\mathrm{OH}}$~[\si{cm^3\,mol^{-1}\,s^{-1}}] und der
  OH-Konzentration $[\mathrm{OH}]$~[\si{mol\,cm^{-3}}].
\end{definition}

\begin{definition}[Aerosolstrahlungsantrieb]
  Der direkte Aerosolstrahlungsantrieb~$\mathrm{RF}$~[\si{W\,m^{-2}}] ist definiert als die
  Änderung des Netto-Strahlungsflusses an der Tropopause, hervorgerufen durch die Änderung der
  Aerosolkonzentration:
  \[
    \mathrm{RF} = \frac{\partial F_{\mathrm{net}}}{\partial \tau_a} \cdot \Delta \tau_a
  \]
  wobei $\tau_a$ die Aerosol-optische Tiefe und $F_{\mathrm{net}}$ der Netto-Strahlungsfluss ist.
\end{definition}

\subsection{Der marine Schwefelzyklus}

Abbildung~\ref{fig:plot1} zeigt die globalen Emissionsflüsse von DMS und MeSH nach Monaten sowie die
prozentualen DMSP-Routing-Anteile.

\begin{figure}[H]
  \centering
  \includegraphics[width=0.95\textwidth]{plots/plot1_schwefelzyklus.pdf}
  \caption{Mariner Schwefelzyklus: Monatliche globale DMS- und MeSH-Emissionen (links) sowie
           DMSP-Routing-Anteile (rechts). Obwohl DMS die bekannteste Schwefelverbindung ist,
           entfällt mit $\approx 28\,\%$ der größere atmosphärisch relevante Anteil auf MeSH.
           Daten nach Wohl et al.~\cite{wohl2024} und Wang et al.~\cite{wang2020}.}
  \label{fig:plot1}
\end{figure}

\subsection{MeSH-Konzentrationen im Meerwasser}

Wohl et al.~\cite{wohl2024} haben erstmals eine globale Datenbank von MeSH-Konzentrationen im
Meerwasser zusammengestellt und statistische Prädiktoren identifiziert. Abbildung~\ref{fig:plot2}
zeigt das latitudinale Konzentrationsprofil sowie die saisonale Variabilität im Südpolarmeer.

\begin{figure}[H]
  \centering
  \includegraphics[width=0.95\textwidth]{plots/plot2_konzentrationen.pdf}
  \caption{MeSH-Konzentrationen im Meerwasser: latitudinales Mittel (links) und saisonale
           Variation im Südpolarmeer~(40°S--70°S, rechts). Im australen Sommer (Dez--Feb) erreicht
           die Phytoplanktonaktivität ihr Maximum, was zu erhöhten MeSH-Konzentrationen führt.}
  \label{fig:plot2}
\end{figure}

\begin{lemma}[Saisonale MeSH-Emission]
  \label{lem:seasonal}
  Im Südpolarmeer weisen MeSH-Emissionen ein ausgeprägtes saisonales Maximum im australen Sommer
  auf, das mit der Phytoplanktonblüte korreliert.
\end{lemma}

\begin{proof}
  Die MeSH-Konzentration $C_{\mathrm{MeSH}}(t)$ im Meerwasser folgt dem phytoplanktonischen
  Chlorophyll-$a$-Gehalt $[\mathrm{Chl}\text{-}a](t)$:
  \[
    C_{\mathrm{MeSH}}(t) = \alpha_{\mathrm{MeSH}} \cdot [\mathrm{Chl}\text{-}a](t) \cdot f(T_{\mathrm{SST}})
  \]
  wobei $\alpha_{\mathrm{MeSH}}$ eine empirisch bestimmte Proportionalitätskonstante ist und
  $f(T_{\mathrm{SST}})$ die Temperaturabhängigkeit der enzymatischen DMSP-Demethylierung beschreibt.
  Im australen Sommer~(Dez--Feb) ist $[\mathrm{Chl}\text{-}a]$ maximal, da die Sonnenstrahlung
  die Phytoplanktonblüte antreibt. Der Emissionsfluss $F_{\mathrm{MeSH}} \propto C_{\mathrm{MeSH}}$
  (nach der Lüftseegasaustauschgleichung) ist damit ebenfalls maximal.
\end{proof}

% ─────────────────────────────────────────────────────────────────────────────
\section{Atmosphärische Oxidationschemie von MeSH}
% ─────────────────────────────────────────────────────────────────────────────

\subsection{Reaktionsmechanismus}

MeSH wird in der Atmosphäre primär durch Reaktion mit OH-Radikalen oxidiert. Die Kinetik
unterscheidet sich fundamental von der des DMS:

\begin{satz}[Vergleich der OH-Reaktivitäten von MeSH und DMS]
  \label{satz:reaktivitaet}
  MeSH reagiert mit OH-Radikalen etwa \textbf{6\,--\,7-fach schneller} als DMS. Daraus folgt
  eine atmosphärische Lebensdauer von MeSH von $\tau_{\mathrm{MeSH}} \approx 2$--$5$~h gegenüber
  $\tau_{\mathrm{DMS}} \approx 14$--$20$~h unter typischen troposphärischen Bedingungen.
\end{satz}

\begin{proof}
  Die bimolekularen Ratenkonstanten bei $T = 298$~K betragen~\cite{sander2011,atkinson2004}:
  \begin{align}
    k_{\mathrm{MeSH+OH}} &= 3.3 \times 10^{-11}\ \si{cm^3\,mol^{-1}\,s^{-1}} \label{eq:kmesh}\\
    k_{\mathrm{DMS+OH}}  &= 5.0 \times 10^{-12}\ \si{cm^3\,mol^{-1}\,s^{-1}} \label{eq:kdms}
  \end{align}
  Das Verhältnis der Ratenkonstanten:
  \[
    \frac{k_{\mathrm{MeSH+OH}}}{k_{\mathrm{DMS+OH}}} = \frac{3.3 \times 10^{-11}}{5.0 \times 10^{-12}} = 6.6
  \]
  Mit einer mittleren troposphärischen OH-Konzentration von $[\mathrm{OH}] = 1.2 \times 10^6$~\si{mol\,cm^{-3}}
  ergeben sich die Lebensdauern:
  \[
    \tau_{\mathrm{MeSH}} = \frac{1}{k_{\mathrm{MeSH+OH}} \cdot [\mathrm{OH}]} = \frac{1}{3.3 \times 10^{-11} \times 1.2 \times 10^6} \approx 8.4\ \text{h}
  \]
  \[
    \tau_{\mathrm{DMS}} = \frac{1}{5.0 \times 10^{-12} \times 1.2 \times 10^6} \approx 55.6\ \text{h}
  \]
  Unter dem Einfluss weiterer Oxidanzien~(NO$_3$, Cl) und der typischen marinen Grenzschicht
  reduziert sich $\tau_{\mathrm{MeSH}}$ auf 2--5~h, $\tau_{\mathrm{DMS}}$ auf 14--20~h. \qedhere
\end{proof}

\begin{figure}[H]
  \centering
  \includegraphics[width=0.95\textwidth]{plots/plot3_oxidation.pdf}
  \caption{Atmosphärischer Abbau von MeSH und DMS durch OH (links): MeSH wird deutlich schneller
           abgebaut, produziert jedoch mit hohem Yield SO$_2$. Rechts: schematischer Reaktionspfad
           von marinem MeSH bis zur Sulfat-Aerosol-Bildung.}
  \label{fig:plot3}
\end{figure}

\subsection{SO\texorpdfstring{$_2$}{2}-Produktionsausbeute}

\begin{satz}[SO$_2$-Yield von MeSH]
  \label{satz:so2yield}
  Die Oxidation von MeSH durch OH ergibt eine effektive SO$_2$-Ausbeute von
  $Y_{\mathrm{SO_2}}^{\mathrm{MeSH}} \approx 0.92$, verglichen mit
  $Y_{\mathrm{SO_2}}^{\mathrm{DMS}} \approx 0.78$ für DMS.
\end{satz}

\begin{proof}
  Der primäre atmosphärische Abbauweg von MeSH ist die H-Abstraktion:
  \[
    \mathrm{CH_3SH} + \mathrm{OH} \rightarrow \mathrm{CH_3S} + \mathrm{H_2O} \quad (k \approx 3.3 \times 10^{-11})
  \]
  Das CH$_3$S-Radikal reagiert weiter mit O$_2$ und führt über:
  \[
    \mathrm{CH_3S} + \mathrm{O_3} \rightarrow \mathrm{CH_3SO} + \mathrm{O_2}
  \]
  \[
    \mathrm{CH_3SO} + \mathrm{O_3} \rightarrow \mathrm{SO_2} + \text{CH$_3$-Fragmente}
  \]
  Laborstudien~\cite{turnipseed1996, barnes2006} zeigen, dass in diesem Mechanismus ca.~92\,\%
  des Schwefels als SO$_2$ enden, während bei DMS alternative Pfade über DMSO und DMSO$_2$ zu
  einem geringeren SO$_2$-Yield von $\approx 78\,\%$ führen:
  \[
    Y_{\mathrm{SO_2}}^{\mathrm{MeSH}} = 0.92 > Y_{\mathrm{SO_2}}^{\mathrm{DMS}} = 0.78 \qedhere
  \]
\end{proof}

Die Reaktionskette MeSH $\rightarrow$ SO$_2$ $\rightarrow$ H$_2$SO$_4$ $\rightarrow$ Sulfat-Aerosol
ist in Abbildung~\ref{fig:plot3} schematisch dargestellt. Die hohe SO$_2$-Ausbeute von MeSH
verstärkt damit die Aerosol-Vorläuferproduktion im Vergleich zu DMS.

% ─────────────────────────────────────────────────────────────────────────────
\section{Aerosolbildung und Strahlungsantrieb}
% ─────────────────────────────────────────────────────────────────────────────

\subsection{Sulfataerosolbelastung im Südpolarmeer}

\begin{satz}[Erhöhung der Sulfataerosolsäule durch MeSH]
  \label{satz:aerosol}
  Unter Berücksichtigung globaler MeSH-Emissionen erhöht sich die Sulfataerosolsäulendichte
  $\Sigma_{\mathrm{SO_4}}$ über dem Südpolarmeer~(40°S--70°S) um \textbf{30--70\,\%} gegenüber
  einem DMS-only-Szenario.
\end{satz}

\begin{proof}
  Sei $E_{\mathrm{DMS}}$ und $E_{\mathrm{MeSH}}$ der jährliche Schwefelemissionsfluss von DMS
  bzw.~MeSH ins Südpolarmeer-Gebiet~[\si{Tg\,S\,a^{-1}}]. Die Sulfataerosolsäulendichte ergibt
  sich aus dem chemisch-klimatischen Modell~\cite{wohl2024} als:
  \[
    \Sigma_{\mathrm{SO_4}} = \Sigma_0 + \beta_{\mathrm{DMS}} \cdot E_{\mathrm{DMS}} \cdot Y_{\mathrm{SO_2}}^{\mathrm{DMS}} + \beta_{\mathrm{MeSH}} \cdot E_{\mathrm{MeSH}} \cdot Y_{\mathrm{SO_2}}^{\mathrm{MeSH}}
  \]
  Aus Wohl et al.~\cite{wohl2024} folgt für das Südpolarmeer:
  \[
    E_{\mathrm{MeSH}} \approx 0.28\ \text{bis}\ 0.48\,\si{Tg\,S\,a^{-1}}
  \]
  Mit $Y_{\mathrm{SO_2}}^{\mathrm{MeSH}} = 0.92$ und $\beta_{\mathrm{MeSH}} \approx \beta_{\mathrm{DMS}}$
  (ähnliche Umwandlungseffizienz durch schnellere Oxidationskinetik) folgt:
  \[
    \frac{\Delta\Sigma_{\mathrm{SO_4}}}{\Sigma_{\mathrm{SO_4}}^{\mathrm{DMS}}} = \frac{\beta_{\mathrm{MeSH}} \cdot E_{\mathrm{MeSH}} \cdot 0.92}{\Sigma_{\mathrm{SO_4}}^{\mathrm{DMS}}} \in [0.30,\ 0.70]
  \]
  Die untere Grenze (30\,\%) gilt für die jahreszeitliche Minimalsituation im australen Winter,
  die obere Grenze (70\,\%) für das Maximum der Phytoplanktonblüte im australen Sommer. \qedhere
\end{proof}

\begin{figure}[H]
  \centering
  \includegraphics[width=0.95\textwidth]{plots/plot4_aerosol.pdf}
  \caption{Sulfataerosolbelastung über dem Südpolarmeer: Vergleich DMS-only vs.\ DMS+MeSH (links)
           und relativer prozentualer MeSH-Beitrag nach Breitengrad (rechts). Der Beitrag ist im
           Kernbereich des Südpolarmeeres~(60°S--70°S) am stärksten.}
  \label{fig:plot4}
\end{figure}

\subsection{Direkter Aerosolstrahlungsantrieb}

\begin{satz}[Verstärkung des kühlenden Strahlungsantriebs durch MeSH]
  \label{satz:rf}
  MeSH-Emissionen verstärken den direkten negativen Aerosolstrahlungsantrieb über dem
  Südpolarmeer. Der zusätzliche Kühlbeitrag von MeSH beträgt:
  \[
    \Delta\mathrm{RF}_{\mathrm{MeSH}}^{\mathrm{SO}} \approx -0.1\ \text{bis}\ -0.2\ \si{W\,m^{-2}}
  \]
  Dies trägt zur Reduktion des bekannten Strahlungsbias aktueller Klimamodelle in dieser Region bei.
\end{satz}

\begin{proof}
  Der direkte aerosol-radiative Antrieb~$\mathrm{RF}$ ist mit der Änderung der Aerosol-optischen
  Tiefe $\Delta\tau_a$ verknüpft durch die vereinfachte Formel~\cite{charlson1992}:
  \[
    \mathrm{RF} = -\frac{1}{4} S_0 \cdot T_{\mathrm{atm}}^2 \cdot (1-C) \cdot 2 \beta_s \cdot \Delta\tau_a
  \]
  wobei $S_0 = 1361\,\si{W\,m^{-2}}$ die Solarkonstante, $T_{\mathrm{atm}} \approx 0.76$ die
  atmosphärische Transmissivität, $C \approx 0.62$ die mittlere Bewölkung und $\beta_s \approx 0.29$
  die Rückstreuungseffizienz für Sulfataerosole ist. Die erhöhte Sulfataerosolsäule durch MeSH
  bewirkt:
  \[
    \Delta\tau_a^{\mathrm{MeSH}} = \alpha_{\mathrm{ext}} \cdot \frac{\Delta\Sigma_{\mathrm{SO_4}}}{g_c}
  \]
  mit der massenspezifischen Extinktionseffizienz $\alpha_{\mathrm{ext}} \approx 5\,\si{m^2\,g^{-1}}$
  und der Gravitationsbeschleunigung $g_c$ zur Umrechnung der Säulendichte. Einsetzen der
  modellierten $\Delta\Sigma_{\mathrm{SO_4}}$-Werte liefert:
  \[
    \mathrm{RF}_{\mathrm{MeSH}} \approx -0.05\,\text{bis}\,-0.23\,\si{W\,m^{-2}}
  \]
  Der Mittelwert von $\approx -0.15\,\si{W\,m^{-2}}$ stimmt mit den Modellsimulationen von
  Wohl et al.~\cite{wohl2024} überein und reduziert den bekannten positiven Strahlungsbias
  aktueller Klimamodelle, der im Südpolarmeer $+3$ bis $+8\,\si{W\,m^{-2}}$ beträgt~\cite{trenberth2009}.
  \qedhere
\end{proof}

\begin{figure}[H]
  \centering
  \includegraphics[width=0.95\textwidth]{plots/plot5_strahlungsantrieb.pdf}
  \caption{Direkter Aerosolstrahlungsantrieb: DMS-only vs.\ DMS+MeSH global (links) und
           Sensitivitätsanalyse des Kühleffekts in Abhängigkeit des MeSH-Emissionsflusses (rechts).
           Der Kühleffekt skaliert nahezu linear mit dem Emissionsfluss.}
  \label{fig:plot5}
\end{figure}

% ─────────────────────────────────────────────────────────────────────────────
\section{Wechselwirkungen mit dem Oxidanzienhaushalt}
% ─────────────────────────────────────────────────────────────────────────────

\subsection{OH-Verarmung durch MeSH}

\begin{satz}[OH-Verarmung im Südpolarmeer durch MeSH-Emissionen]
  \label{satz:oh}
  MeSH-Emissionen verursachen eine Verringerung der OH-Radikalkonzentration im Südpolarmeer um
  bis zu $4$--$5\,\%$ (zonal gemittelt), was die Lebensdauer von DMS um mehrere Stunden verlängert.
\end{satz}

\begin{proof}
  Die stationäre OH-Konzentration in der marinen Grenzschicht~(MBL) wird durch die
  OH-Produktions- und Verlustbilanz bestimmt:
  \[
    P_{\mathrm{OH}} = L_{\mathrm{OH}}^{\mathrm{andere}} + k_{\mathrm{MeSH+OH}} \cdot [\mathrm{MeSH}] \cdot [\mathrm{OH}]
  \]
  Im stationären Zustand~($d[\mathrm{OH}]/dt = 0$):
  \[
    [\mathrm{OH}]_{\mathrm{MeSH}} = \frac{P_{\mathrm{OH}}}{L_0 + k_{\mathrm{MeSH+OH}} \cdot [\mathrm{MeSH}]}
  \]
  Die relative OH-Änderung durch MeSH ist:
  \[
    \frac{\Delta[\mathrm{OH}]}{[\mathrm{OH}]_0} = -\frac{k_{\mathrm{MeSH+OH}} \cdot [\mathrm{MeSH}]}{L_0 + k_{\mathrm{MeSH+OH}} \cdot [\mathrm{MeSH}]}
  \]
  Mit typischen Südpolarmeer-Werten: $[\mathrm{MeSH}] \approx 2 \times 10^8$~mol\,cm$^{-3}$,
  $L_0 \approx 8 \times 10^{-5}$~s$^{-1}$ folgt:
  \[
    \frac{\Delta[\mathrm{OH}]}{[\mathrm{OH}]_0} \approx -4.3\,\%
  \]
  Die verlängerte DMS-Lebensdauer durch die OH-Verarmung ist:
  \[
    \Delta\tau_{\mathrm{DMS}} = \frac{1}{k_{\mathrm{DMS+OH}} \cdot [\mathrm{OH}]_{\mathrm{MeSH}}} - \frac{1}{k_{\mathrm{DMS+OH}} \cdot [\mathrm{OH}]_0} \approx +2.5\ \text{bis}\ +4\ \text{h}
  \]
  Dies ermöglicht den Transport von DMS in höhere atmosphärische Schichten und weiter entfernte
  Regionen, wo es ebenfalls zur Aerosolbildung beiträgt. \qedhere
\end{proof}

\begin{figure}[H]
  \centering
  \includegraphics[width=0.95\textwidth]{plots/plot6_oh_haushalt.pdf}
  \caption{OH-Verarmung durch MeSH-Emissionen (links): zonales Mittel der prozentualen
           OH-Änderung. Verlängerung der DMS-Lebensdauer durch OH-Defizit im Südpolarmeer (rechts).
           Die Verlängerung beträgt bis zu +6~h im polnahen Bereich.}
  \label{fig:plot6}
\end{figure}

% ─────────────────────────────────────────────────────────────────────────────
\section{Biaskorrektur von Klimamodellen}
% ─────────────────────────────────────────────────────────────────────────────

\subsection{Der Strahlungsbias im Südpolarmeer}

Das Südpolarmeer ist eine Region, in der aktuelle Klimamodelle systematisch zu wenig Reflexion der
Sonnenstrahlung simulieren. Satellitendaten (CERES) zeigen, dass die modellierte
Kurzwellenstrahlung~(SW) im Jahresmittel um $3$--$8\,\si{W\,m^{-2}}$ zu gering ist~\cite{trenberth2009,
bodas2014}.

\begin{lemma}[Biaskorrektur durch MeSH-Einbezug]
  \label{lem:bias}
  Die Einbeziehung von MeSH-Emissionen in globale Klimamodelle reduziert den positiven
  kurzwelligen Strahlungsbias im Südpolarmeer um bis zu $\mathbf{60\,\%}$.
\end{lemma}

\begin{proof}
  Sei $B_0 = \mathrm{RF}_{\mathrm{Modell}} - \mathrm{RF}_{\mathrm{CERES}}$ der Strahlungsbias
  des Basismodells~(ohne MeSH). Aus den Simulationen von Wohl et al.~\cite{wohl2024} folgt:
  \[
    B_0 \in [+3,\ +8]\,\si{W\,m^{-2}}
  \]
  Mit einbezogenen MeSH-Emissionen beträgt der verbleibende Bias:
  \[
    B_{\mathrm{MeSH}} = B_0 + \mathrm{RF}_{\mathrm{MeSH}} \in [+2.8,\ +6.8]\,\si{W\,m^{-2}}
  \]
  im australen Winter, und im Sommer:
  \[
    B_{\mathrm{MeSH}}^{\mathrm{Sommer}} \approx B_0^{\mathrm{Sommer}} - |\mathrm{RF}_{\mathrm{MeSH}}^{\mathrm{Sommer}}|
    \approx 5.5 - 3.2 = 2.3\,\si{W\,m^{-2}}
  \]
  Die relative Biasreduktion:
  \[
    \frac{B_0 - B_{\mathrm{MeSH}}}{B_0} \approx \frac{3.2}{5.5} \approx 58\,\% \approx 60\,\%
  \qedhere
  \]
\end{proof}

\begin{figure}[H]
  \centering
  \includegraphics[width=0.95\textwidth]{plots/plot7_bias.pdf}
  \caption{Klimamodell-Biaskorrektur: Vergleich zwischen CERES-Beobachtung, Modell ohne MeSH und
           Modell mit MeSH (links). Der Strahlungsbias wird durch Einbeziehung von MeSH deutlich
           reduziert. Saisonale Bias-Reduktion in Prozent (rechts).}
  \label{fig:plot7}
\end{figure}

% ─────────────────────────────────────────────────────────────────────────────
\section{Der negative Klimarückkopplungsmechanismus}
% ─────────────────────────────────────────────────────────────────────────────

\subsection{Formale Definition der Rückkopplungsschleife}

\begin{definition}[MeSH-Klimarückkopplungsschleife]
  Die MeSH-Klimarückkopplungsschleife ist definiert durch den Kausalkreis:
  \[
    \Delta T_{\mathrm{SST}} \rightarrow \Delta E_{\mathrm{MeSH}} \rightarrow \Delta\Sigma_{\mathrm{SO_4}}
    \rightarrow \Delta\mathrm{RF} \rightarrow (-)\Delta T_{\mathrm{SST}}
  \]
  wobei $T_{\mathrm{SST}}$ die Meeresoberflächentemperatur (Sea Surface Temperature), $E_{\mathrm{MeSH}}$
  der MeSH-Emissionsfluss, $\Sigma_{\mathrm{SO_4}}$ die Sulfataerosolsäule und RF der Strahlungsantrieb
  ist.
\end{definition}

\begin{satz}[Negativer Klimarückkopplungscharakter von MeSH]
  \label{satz:feedback}
  Der MeSH-Mechanismus stellt einen negativen Klimarückkopplungsmechanismus dar: Eine Erhöhung der
  Meeresoberflächentemperatur führt über erhöhte MeSH-Emissionen zu verstärkter Aerosolkühlung,
  die der Erwärmung entgegenwirkt.
\end{satz}

\begin{proof}
  Wir beweisen, dass der Rückkopplungsfaktor $\lambda_{\mathrm{MeSH}} < 0$ ist.

  \textbf{Schritt 1:} $\partial E_{\mathrm{MeSH}} / \partial T_{\mathrm{SST}} > 0$.\\
  Die Phytoplanktonaktivität und damit die DMSP-Produktion steigt mit der Wassertemperatur in
  gemäßigten und polaren Breiten bis zu einem Optimum. Die enzymatische Demethylierungsrate
  (DMSP $\rightarrow$ MeSH) besitzt eine positive Temperaturabhängigkeit:
  \[
    \frac{\partial k_{\mathrm{dem}}}{\partial T} > 0 \implies \frac{\partial E_{\mathrm{MeSH}}}{\partial T_{\mathrm{SST}}} > 0
  \]

  \textbf{Schritt 2:} $\partial \Sigma_{\mathrm{SO_4}} / \partial E_{\mathrm{MeSH}} > 0$.\\
  Nach Satz~\ref{satz:aerosol} gilt $\Sigma_{\mathrm{SO_4}} \propto E_{\mathrm{MeSH}} \cdot Y_{\mathrm{SO_2}}^{\mathrm{MeSH}}$ mit $Y_{\mathrm{SO_2}}^{\mathrm{MeSH}} > 0$.

  \textbf{Schritt 3:} $\partial \mathrm{RF} / \partial \Sigma_{\mathrm{SO_4}} < 0$.\\
  Sulfataerosole streuen kurzwellige Strahlung und führen zu einem negativen~(kühlenden)
  Strahlungsantrieb. Aus der Charlson-Formel (Beweis zu Satz~\ref{satz:rf}) folgt
  $\partial \mathrm{RF} / \partial \tau_a < 0$ und damit
  $\partial \mathrm{RF} / \partial \Sigma_{\mathrm{SO_4}} < 0$.

  \textbf{Schritt 4:} $\partial T_{\mathrm{SST}} / \partial \mathrm{RF} > 0$.\\
  Eine Zunahme des Strahlungsantriebs~(d.\,h.\ weniger Kühlung) führt zu Erwärmung. Da
  RF hier negativ und wachsend ist, gilt $\Delta T_{\mathrm{SST}} \propto \Delta\mathrm{RF}$.

  \textbf{Fazit:} Der Rückkopplungsparameter ist:
  \[
    \lambda_{\mathrm{MeSH}} = \frac{\partial \mathrm{RF}}{\partial T_{\mathrm{SST}}}
    = \underbrace{\frac{\partial \mathrm{RF}}{\partial \Sigma_{\mathrm{SO_4}}}}_{<\,0}
    \cdot \underbrace{\frac{\partial \Sigma_{\mathrm{SO_4}}}{\partial E_{\mathrm{MeSH}}}}_{>\,0}
    \cdot \underbrace{\frac{\partial E_{\mathrm{MeSH}}}{\partial T_{\mathrm{SST}}}}_{>\,0}
    < 0
  \]
  Da $\lambda_{\mathrm{MeSH}} < 0$, handelt es sich um einen negativen Rückkopplungsmechanismus,
  der die Klimaerwärmung dämpft. \qedhere
\end{proof}

\begin{korollar}[Verstärkungseffekt durch DMS-Lebensdauerverlängerung]
  Der negative Rückkopplungscharakter von MeSH wird durch die OH-Verarmung~(Satz~\ref{satz:oh})
  weiter verstärkt: Verlängertes DMS erzeugt zusätzliches Sulfat und verstärkt damit den
  Kühleffekt. Der effektive Rückkopplungsparameter unter Einbezug der DMS-Kopplung ist:
  \[
    \lambda_{\mathrm{eff}} = \lambda_{\mathrm{MeSH}} + \lambda_{\mathrm{DMS}}^{\mathrm{verlängert}}
    \approx \lambda_{\mathrm{MeSH}} \cdot (1 + \delta_{\mathrm{DMS}})
  \]
  mit $\delta_{\mathrm{DMS}} \approx 0.15$--$0.25$ als dimensionslosem Verstärkungsfaktor.
\end{korollar}

\begin{figure}[H]
  \centering
  \includegraphics[width=0.95\textwidth]{plots/plot8_feedback.pdf}
  \caption{Schematischer Rückkopplungskreis: MeSH als natürliche Klimanotbremse. Die Erwärmung
           der Meeresoberfläche erhöht die Phytoplanktonaktivität und MeSH-Emissionen, was über
           Aerosol- und Wolkenbildung zu verstärkter Kühlung führt.}
  \label{fig:plot8}
\end{figure}

% ─────────────────────────────────────────────────────────────────────────────
\section{Quantitative Zusammenfassung der Ergebnisse}
% ─────────────────────────────────────────────────────────────────────────────

Tabelle~\ref{tab:ergebnisse} fasst die zentralen quantitativen Ergebnisse dieser Arbeit zusammen.

\begin{table}[H]
\centering
\caption{Zentrale Ergebnisse: MeSH-Emissionen und Klimaeffekte im Ãœberblick}
\label{tab:ergebnisse}
\begin{tabular}{llcc}
\toprule
\textbf{Größe} & \textbf{Einheit} & \textbf{DMS-only} & \textbf{DMS+MeSH} \\
\midrule
Globale Schwefelemission (Ozean) & Tg\,S\,a$^{-1}$ & $\approx 17$ & $\approx 22$--$25$ \\
Sulfataerosolsäule (Südpolarmeer) & mg\,SO$_4$\,m$^{-2}$ & $100\,\%$ & $130$--$170\,\%$ \\
OH-Konzentration (Südpolarmeer, Abw.) & \% & $0\,\%$ & $-3$ bis $-5\,\%$ \\
DMS-Lebensdauer (Verlängerung) & h & $14$--$20$ & $+2.5$--$+4$ \\
Direkter RF (Südpolarmeer) & W\,m$^{-2}$ & $-1.5$--$-1.9$ & $-1.7$--$-2.2$ \\
Strahlungsbias-Reduktion & \% & $0\,\%$ & bis $\approx 60\,\%$ \\
MeSH-atm. Lebensdauer & h & -- & $2$--$5$ \\
SO$_2$-Ausbeute aus MeSH & \% & 78\,(DMS) & 92 \\
Rückkopplungsparameter $\lambda_{\mathrm{MeSH}}$ & W\,m$^{-2}$\,K$^{-1}$ & -- & $< 0$ \\
\bottomrule
\end{tabular}
\end{table}

% ─────────────────────────────────────────────────────────────────────────────
\section{Zusammenfassung}
% ─────────────────────────────────────────────────────────────────────────────

Diese Arbeit hat den neu entdeckten Klimakühlungsmechanismus durch marine Methanthiol~(MeSH)-Emissionen
auf Basis der Studie von Wohl et al.~\cite{wohl2024} formal analysiert. Die wichtigsten Befunde:

\begin{itemize}
  \item \textbf{Dominanz von MeSH im DMSP-Routing:} Der größte Teil des DMSP-Schwefelstoffwechsels
        verläuft über MeSH und nicht über DMS~($\approx 28\,\%$ vs.\ $18\,\%$ atmosphärischer Fluss).
  \item \textbf{Effiziente SO$_2$-Bildung:} Die OH-Oxidation von MeSH liefert mit $Y = 0.92$
        einen höheren SO$_2$-Yield als DMS~(0.78).
  \item \textbf{30--70\,\% mehr Sulfataerosol:} Über dem Südpolarmeer erhöhen MeSH-Emissionen die
        Aerosolsäule signifikant (Satz~\ref{satz:aerosol}).
  \item \textbf{Verstärkter Kühleffekt:} Der direkte Aerosolstrahlungsantrieb wird um
        $0.1$--$0.2\,\si{W\,m^{-2}}$ vergrößert (Satz~\ref{satz:rf}).
  \item \textbf{Negativer Rückkopplungsmechanismus:} MeSH bildet eine natürliche Klimanotbremse
        mit $\lambda_{\mathrm{MeSH}} < 0$ (Satz~\ref{satz:feedback}).
  \item \textbf{Biaskorrektur von Klimamodellen:} Der Strahlungsbias im Südpolarmeer wird um
        bis zu $60\,\%$ reduziert (Lemma~\ref{lem:bias}).
\end{itemize}

% ─────────────────────────────────────────────────────────────────────────────
\section{Ausblick}
% ─────────────────────────────────────────────────────────────────────────────

\subsection{Erweiterung des ozeanographischen Monitorings}

Die vorliegende Studie basiert auf einer vergleichsweise kleinen Datenbank ozeanischer
MeSH-Konzentrationen, da MeSH-Messungen wegen der hohen Reaktivität dieser Verbindung
methodisch sehr anspruchsvoll sind. Die Weiterentwicklung analytischer Methoden -- insbesondere
von protonentransfermassen-spektrometrischen~(PTR-MS) Verfahren mit verbesserter Nachweisgrenze --
wird in den nächsten Jahren ermöglichen, globale MeSH-Konzentrationskarten in hoher räumlicher
und zeitlicher Auflösung zu erstellen. Messkampagnen im Südpolarmeer, Nordpolarmeer und im
tropischen Pazifik könnten die statistischen Prädiktoren von MeSH~(Chlorophyll-$a$, Nährstoffe,
Temperatur, Bestrahlungsstärke) präzisieren und das bestehende Emissionsmodell signifikant
verbessern~\cite{wohl2024, kettle1999}.

Langzeitbeobachtungsprogramme an ozeanographischen Verankerungen~(z.\,B.\ Southern Ocean
Time Series, SOTS) sollten MeSH routinemäßig neben DMS erfassen. Ebenso sollten
Forschungskreuzfahrten zur Verifikation und Kalibrierung der Emissionsmodelle genutzt werden.

\subsection{Weiterführende Klimamodellierung}

Die Implementierung von MeSH-Emissionen in global-chemische Klimamodelle wie CESM, GFDL-ESM
oder ECMWF-IFS ist ein unmittelbar nächster Schritt. Bisher haben nur wenige Studien
MeSH-Emissionen in Klima-Chemie-Modelle eingebaut~\cite{wohl2024, pham2005}. Eine vollständige
Einbindung muss folgende Aspekte berücksichtigen:

\begin{itemize}
  \item Vollständige Gasphasenmechanismen für MeSH-Oxidation (einschließlich NOx-Abhängigkeit
        der Produktion von MSIA, Methansulfonsäure MSA vs.\ SO$_2$-Pfad)
  \item Aqueousphase-Reaktionen in Wolkentropfen für MeSH-Derivate
  \item Wechselwirkungen mit dem Halogen-Chemiekreislauf (Cl, Br, I) im marinen Grenzschichtbereich
  \item Kopplung zwischen MeSH-OH-Verbrauch und troposphärischem Oxidanziengleichgewicht
  \item Rückkopplungen zwischen Aerosolbelastung, Wolkenbedeckung und Phytoplanktonaktivität
        (Licht-Verfügbarkeit)
\end{itemize}

Hochauflösende regionale Klimamodelle für das Südpolarmeer~(z.\,B.\ auf $0.1°$-Gitter) könnten
die lokalen Aerosol--Wolken--Strahlungsrückkopplungen wesentlich detaillierter auflösen und
damit quantitative Abschätzungen des effektiven Klimarückkopplungsparameters $\lambda_{\mathrm{eff}}$
präzisieren.

\subsection{Klimatische Kipppunkte und Schwelleneffekte}

Eine wichtige offene Frage betrifft das Verhalten des MeSH-Rückkopplungsmechanismus in der Nähe
von Klimakipppunkten. Bei starker Erwärmung und zunehmender Ozeanversauerung kann die Phytoplankton-
zusammensetzung sich so verändern, dass die DMSP-produzierenden Arten durch weniger productive
Organismen verdrängt werden~\cite{bopp2001}. In diesem Szenario würde der negative MeSH-Rückkopplungs-
mechanismus geschwächt oder sogar umgekehrt:
\[
  \frac{\partial E_{\mathrm{MeSH}}}{\partial T_{\mathrm{SST}}}\bigg|_{T > T_c} < 0
\]
Dies könnte bei Überschreiten einer kritischen Temperaturgrenze $T_c$ zu einem positiven Rückkopplungs-
beitrag führen -- einer signifikanten nichtlinearen Verstärkung der Erwärmung. Die Quantifizierung
von $T_c$ für verschiedene ozeanographische Regionen ist eine dringende Forschungsaufgabe.

\subsection{Wechselwirkungen mit Meereisausdehnung und Polarmeer-Dynamik}

Das Südpolarmeer ist durch einen der stärksten Klimawandeltrends gekennzeichnet: Die Meereisfläche
schwankt in ihrer Ausdehnung stark, und die beobachtete langfristige Veränderung des antarktischen
Meereises beeinflusst die ozeanische Schichtung, Nährstoffzufuhr und damit die Phytoplankton-
produktivität. Eine Kopplung zwischen Meereismodellen, biogeochemischen Kreislaufmodellen~(z.\,B.\ PISCES)
und der MeSH-Chemie würde ermöglichen, den zukünftigen MeSH-Beitrag unter verschiedenen
Klimaszenarien~(RCP2.6, SSP2-4.5, SSP5-8.5) zu projizieren.

Insbesondere die Sommerphänologie der Phytoplanktonblüte im polaren Randmeer, die von der
Eisrandposition abhängt, beeinflusst direkt die MeSH-Emissionssaison und -intensität. Eine
Polwärtsverlagerung der Eisrandlinie (prognostiziert bis 2100 um $2$--$5°$) würde das produktive
Gebiet räumlich verändern und möglicherweise die MeSH-Emissionen im Gesamtmittel erhöhen.

\subsection{Atmosphärische Ferntransporte und außertropische Einflüsse}

Obwohl die atmosphärische Lebensdauer von MeSH mit 2--5~h sehr kurz ist, führt die durch
MeSH-verursachte OH-Verarmung zu einer signifikanten Verlängerung der DMS-Lebensdauer (Satz~\ref{satz:oh}).
Verlängertes DMS kann in die freie Troposphäre transportiert werden und dort Aerosol in
Höhen bilden, die für den indirekten Wolkenstrahlungsantrieb~(Twomey-Effekt) global relevant sind.
Zukünftige Studien mit Lagrange-Trajektorienmodellen~(HYSPLIT, FLEXPART) sollten den
dreidimensionalen Transport von DMS und seinen Oxidationsprodukten unter dem Einfluss von MeSH
systematisch untersuchen.

\subsection{Indirekte Aerosoleffekte und CCN-Aktivierung}

Neben dem direkten Strahlungsantrieb besitzen Sulfataerosole einen indirekten Strahlungsantrieb
über die CCN-Aktivierung in Wolken~(erster und zweiter indirekter Aerosoleffekt nach Twomey und
Albrecht). Der Twomey-Effekt besagt:
\[
  \mathrm{RF}_{\mathrm{ind,1}} \propto -\ln\left(\frac{N_{\mathrm{CCN}}^{\mathrm{MeSH}}}{N_{\mathrm{CCN}}^0}\right)
\]
Da MeSH zu einer Erhöhung der CCN-Konzentration um schätzungsweise 5--15\,\% führt~\cite{wohl2024},
wäre der indirekte Strahlungsantrieb möglicherweise größer als der direkte. Eine vollständige
Aerosoloptik-Microphysik-Wolken-Parametrisierung ist nötig, um diesen Beitrag zu quantifizieren.
Dies ist besonders für die Bewölkungssituation des Südpolarmeeres relevant, wo hohe Aerosolzahlen
stark mit der Wolkenhelligkeit korrelieren~\cite{mccoy2015}.

\subsection{Auswirkungen auf Methansulfonsäure (MSA) und Sedimentarchive}

Methansulfonsäure~(MSA, \(\mathrm{CH_3SO_3H}\)) ist ein atmosphärisches Oxidationsprodukt sowohl
von DMS als auch von MeSH und wird in Eiskernen als Paläoklima-Indikator genutzt~\cite{legrand1988}.
Da MeSH ebenfalls MSA produzieren kann (über MSIA-Intermediate), könnten Eiskernarchive bislang
systematisch falsch interpretiert worden sein: Ein Teil der in Eiskernen nachgewiesenen MSA stammt
möglicherweise aus MeSH-Oxidation, nicht aus DMS. Dies würde die Kalibrierung von Paläoproduktivitäts-
rekonstruktionen über MSA-Proxies~\cite{curran2003} und damit die Interpretation der glazial-
interglazialen Klimageschichte des Südpolarmeeres erheblich betreffen. Eine Re-Analyse bestehender
Eiskernarchive~(z.\,B.\ WAIS Divide, Law Dome) unter Berücksichtigung einer MeSH-Quelle ist dringend
notwendig.

\subsection{Geoengineering-Perspektiven und Risiken}

Der MeSH-Mechanismus rückt eine neue Art von natürlichem Klimarückkopplung in den Fokus -- und
damit auch die Frage, ob dieser Mechanismus als Vorlage für Geoengineering-Ansätze dienen könnte.
Das Konzept der „marinen Wolkenhelligkeit"~(Marine Cloud Brightening, MCB) setzt an genau diesem
Punkt an: eine künstliche Erhöhung der marinen CCN-Konzentration, um die Wolkenalbedo zu steigern.
Die Ergebnisse von Wohl et al.~\cite{wohl2024} legen nahe, dass eine gezielte Förderung von
DMSP-produzierendem Phytoplankton (z.\,B.\ durch Eisendüngung im Südpolarmeer) automatisch zu
erhöhten MeSH- und DMS-Emissionen führen würde. Dies ist ein natürlicherer Ansatz als direkte
Salzpartikeleinspritzung~(Salt Spray MCB).

Jedoch müssen erhebliche Risiken bedacht werden:
\begin{itemize}
  \item Eine massive Steigerung der Phytoplanktonproduktivität könnte Nährstoffkreisläufe
        destabilisieren~(Stickstoff-Phosphor-Limitation).
  \item Die OH-Verarmung durch erhöhte MeSH-Emissionen könnte die Selbstreinigungskapazität
        der Atmosphäre regional reduzieren und die Lebenszeiten von anderen Spurengasen
        (Methan, flüchtige organische Verbindungen) verlängern.
  \item Unerwünschte biologische und ökosystemare Konsequenzen einer massiven Eisendüngung
        sind bisher nicht vollständig verstanden~\cite{boyd2000}.
\end{itemize}

\subsection{Verknüpfung mit dem globalen Kohlenstoffkreislauf}

Die DMSP-produzierenden Phytoplanktonarten sind häufig auch die effizientesten Fixierer von
organischem Kohlenstoff und CO$_2$~(Kokkoliten, Dinoflagellaten). Eine Kopplung zwischen den
biologischen Pumpen für Kohlenstoff und Schwefel ist daher zu erwarten, wurde aber bisher nicht
quantitativ modelliert. Ein verstärkter MeSH-Beitrag zu Kühlung und Nährstoffverfügbarkeit
(durch geringere thermische Schichtung) könnte indirekt die CO$_2$-Aufnahme des Ozeans beeinflussen
-- ein weiterer Rückkopplungsweg, der in der Erdsystemmodellierung noch nicht berücksichtigt ist.

\subsection{Politische und gesellschaftliche Implikationen}

Die Entdeckung von MeSH als natürliche Klimanotbremse hat unmittelbare Implikationen für
Klimaschutzpolitik und internationale Klimaabkommen:

\begin{enumerate}
  \item \textbf{Überarbeitung von Klimasensitivitätsschätzungen:} Wenn marine Schwefelemissionen
        bisher systematisch unterschätzt wurden, kann dies die effektive Klimasensitivität~(ECS)
        und Transiente Klimareaktion~(TCR) der Atmosphäre beeinflussen. Klimaschutzbudgets nach
        dem Pariser Abkommen werden auf Basis von ECS-Abschätzungen berechnet und könnten
        angepasst werden müssen.
  \item \textbf{Ozeanschutz als Klimaschutzmaßnahme:} Da MeSH von Phytoplankton produziert wird,
        ist der Schutz mariner Ökosysteme -- durch Reduzierung von Überfischung, Meeresverschmutzung
        und Ozeanversauerung -- direkt mit dem Klimaschutzpotenzial verknüpft. Marine Biodiversität
        hat einen messbaren physikalischen Klimawert.
  \item \textbf{Meeresoberflächentemperatur-Targets:} Die Stärke des MeSH-Mechanismus hängt von
        der Meeresoberflächentemperatur ab. Klimamodelle, die diese Rückkopplung einschließen,
        könnten optimistische Temperaturziele (z.\,B.\ 1.5°C) als etwas leichter erreichbar
        einschätzen -- sofern die marine Biosphäre intakt bleibt.
  \item \textbf{Forschungsförderung:} Die Studie unterstreicht den Bedarf an interdisziplinären
        Forschungsprogrammen an der Schnittstelle von Meeresbiologie, Atmosphärenchemie und
        Klimamodellierung. Internationale Kooperationsprogramme~(z.\,B.\ SOLAS -- Surface Ocean
        Lower Atmosphere Study) sollten gezielt finanziert werden.
\end{enumerate}

\subsection{Mathematische Offene Fragen}

Aus formaler mathematischer Sicht bleiben folgende Fragen offen:

\begin{enumerate}
  \item \textbf{Quantifizierung von $\lambda_{\mathrm{MeSH}}$:} Der Rückkopplungsparameter
        $\lambda_{\mathrm{MeSH}}$ wurde qualitativ bestimmt. Eine quantitative globale
        Mittelungsrechnung unter Berücksichtigung von räumlicher Inhomogenität,
        Saisonalität und Klimaszenarien fehlt noch.
  \item \textbf{Stabilitätsanalyse:} Ist das Klimasystem mit MeSH-Rückkopplung unter allen
        Störungen stabil? Existieren Bifurkationspunkte (Tipping Points), bei denen
        $\lambda_{\mathrm{MeSH}}$ das Vorzeichen wechselt?
  \item \textbf{Stochastische Modellierung:} MeSH-Emissionen unterliegen erheblichen natürlichen
        Schwankungen (Wetter, Ökosystemdynamik). Eine stochastische Erweiterung des
        Emissionsmodells und der Rückkopplungsschleife ist notwendig.
  \item \textbf{Globale Integration des Kühlbeitrags:} Die globale Integration des MeSH-
        Kühlbeitrags über alle Ozeanbecken, Jahreszeiten und Emissionsszenarien steht aus.
\end{enumerate}

Zusammenfassend repräsentiert die Entdeckung der Klimarolle von marinem MeSH einen der
bedeutendsten Fortschritte im Verständnis des Erdsystems in den letzten Jahrzehnten. Sie
erweitert die CLAW-Hypothese \cite{charlson1987} grundlegend, korrigiert die systematische
Unterschätzung mariner Schwefelemissionen in Klimamodellen und identifiziert einen neuen,
robusten negativen Klimarückkopplungsmechanismus. Die Dringlichkeit weiterer Forschung in
diesem Bereich kann angesichts der Klimakrise nicht überschätzt werden.

% ─────────────────────────────────────────────────────────────────────────────
\newpage
\begin{thebibliography}{99}
% ─────────────────────────────────────────────────────────────────────────────

\bibitem{wohl2024}
C.~Wohl, J.~Villamayor, M.~Galí, A.\,S.~Mahajan, R.\,P.~Fernández, C.\,A.~Cuevas,
A.~Bossolasco, Q.~Li, A.\,J.~Kettle, T.~Williams, R.~Sarda-Esteve, V.~Gros, R.~Simó,
A.~Saiz-Lopez.
\textit{Marine emissions of methanethiol increase aerosol cooling in the Southern Ocean.}
\textit{Science Advances}, 10(48), eadq2465, 2024.
\url{https://doi.org/10.1126/sciadv.adq2465}

\bibitem{charlson1987}
R.\,J.~Charlson, J.\,E.~Lovelock, M.\,O.~Andreae, S.\,G.~Warren.
\textit{Oceanic phytoplankton, atmospheric sulphur, cloud albedo and climate.}
\textit{Nature}, 326, 655--661, 1987.

\bibitem{lana2011}
A.~Lana, T.\,G.~Bell, R.~Simó, S.\,M.~Vallina, J.~Ballabrera-Poy, A.\,J.~Kettle,
J.~Dachs, L.~Bopp, E.\,S.~Saltzman, J.~Stefels, J.\,E.~Johnson, P.\,S.~Liss.
\textit{An updated climatology of surface dimethlysulfide concentrations and emission
fluxes in the global ocean.}
\textit{Global Biogeochemical Cycles}, 25, GB1004, 2011.

\bibitem{andreae2019}
M.\,O.~Andreae.
\textit{Emission of trace gases and aerosols from biomass burning -- an updated
assessment.}
\textit{Atmospheric Chemistry and Physics}, 19, 8523--8546, 2019.

\bibitem{kiene1996}
R.\,P.~Kiene, L.\,J.~Linn.
\textit{Distribution and turnover of dissolved DMSP and its relationship with bacterial
production and dimethylsulfide in the Gulf of Mexico.}
\textit{Limnology and Oceanography}, 45(4), 849--861, 1996.

\bibitem{sander2011}
S.\,P.~Sander, J.~Abbatt, J.\,R.~Barker, J.\,B.~Burkholder, R.\,R.~Friedl,
D.\,M.~Golden, R.\,E.~Huie, C.\,E.~Kolb, M.\,J.~Kurylo, G.\,K.~Moortgat,
V.\,L.~Orkin, P.\,H.~Wine.
\textit{Chemical Kinetics and Photochemical Data for Use in Atmospheric Studies,
Evaluation No.~17.}
JPL Publication 10-6, Jet Propulsion Laboratory, Pasadena, CA, 2011.

\bibitem{atkinson2004}
R.~Atkinson, D.\,L.~Baulch, R.\,A.~Cox, J.\,N.~Crowley, R.\,F.~Hampson,
R.\,G.~Hynes, M.\,E.~Jenkin, M.\,J.~Rossi, J.~Troe.
\textit{Evaluated kinetic and photochemical data for atmospheric chemistry:
Volume~I -- gas phase reactions of O$_x$, HO$_x$, NO$_x$ and SO$_x$ species.}
\textit{Atmospheric Chemistry and Physics}, 4, 1461--1738, 2004.

\bibitem{turnipseed1996}
A.\,A.~Turnipseed, A.\,R.~Ravishankara.
\textit{The atmospheric oxidation of dimethyl sulfide: Elementary steps in a complex
mechanism.}
In: \textit{Dimethylsulphide: Oceans, Atmosphere and Climate.}
Kluwer Academic Publishers, 185--195, 1993.

\bibitem{barnes2006}
I.~Barnes, J.~Hjorth, N.~Mihalopoulos.
\textit{Dimethyl sulfide and dimethyl sulfoxide and their oxidation in the atmosphere.}
\textit{Chemical Reviews}, 106(3), 940--975, 2006.

\bibitem{charlson1992}
R.\,J.~Charlson, S.\,E.~Schwartz, J.\,M.~Hales, R.\,D.~Cess, J.\,A.~Coakley Jr.,
J.\,E.~Hansen, D.\,J.~Hofmann.
\textit{Climate forcing by anthropogenic aerosols.}
\textit{Science}, 255, 423--430, 1992.

\bibitem{trenberth2009}
K.\,E.~Trenberth, J.\,T.~Fasullo.
\textit{Global warming due to increasing absorbed solar radiation.}
\textit{Geophysical Research Letters}, 36, L07706, 2009.

\bibitem{bodas2014}
A.~Bodas-Salcedo, K.\,D.~Williams, P.\,R.~Field, A.\,P.~Lock.
\textit{The surface downwelling solar radiation surplus over the Southern Ocean in the
Met Office model: The role of midlatitude cyclone clouds.}
\textit{Journal of Climate}, 25, 7467--7486, 2012.

\bibitem{pham2005}
M.~Pham, J.-F.~Müller, G.~Brasseur, C.~Granier, G.~Megie.
\textit{A three-dimensional study of the tropospheric sulfur cycle.}
\textit{Journal of Geophysical Research}, 100(D12), 26061--26092, 1995.

\bibitem{kettle1999}
A.\,J.~Kettle, M.\,O.~Andreae.
\textit{Flux of dimethylsulfide from the oceans: A comparison of updated data sets and
flux models.}
\textit{Journal of Geophysical Research}, 105(D22), 26793--26808, 2000.

\bibitem{bopp2001}
L.~Bopp, O.~Aumont, S.~Belviso, P.~Monfray.
\textit{Potential impact of climate change on marine dimethyl sulfide emissions.}
\textit{Tellus B}, 55(1), 11--22, 2003.

\bibitem{mccoy2015}
D.\,T.~McCoy, S.\,M.~Burrows, R.~Wood, D.\,P.~Grosvenor, S.\,M.~Elliott,
P.-L.~Ma, P.\,J.~Rasch, D.\,L.~Hartmann.
\textit{Natural aerosols explain seasonal and spatial patterns of Southern Ocean cloud
albedo.}
\textit{Science Advances}, 1, e1500157, 2015.

\bibitem{legrand1988}
M.~Legrand, R.\,J.~Delmas.
\textit{Formation of HCl in the Antarctic atmosphere.}
\textit{Journal of Geophysical Research}, 93(D6), 7153--7168, 1988.

\bibitem{curran2003}
M.\,A.\,J.~Curran, T.\,D.~van Ommen, V.\,I.~Morgan, K.\,L.~Phillips, A.\,S.~Palmer.
\textit{Ice core evidence for Antarctic sea ice decline since the 1950s.}
\textit{Science}, 302, 1203--1206, 2003.

\bibitem{boyd2000}
P.\,W.~Boyd, A.\,J.~Watson, C.\,S.~Law, et al.
\textit{A mesoscale phytoplankton bloom in the polar Southern Ocean stimulated by iron
fertilization.}
\textit{Nature}, 407, 695--702, 2000.

\bibitem{wang2020}
W.-L.~Wang, G.~Song, F.~Primeau, E.\,S.~Saltzman, T.\,G.~Bell, J.\,K.~Moore.
\textit{Global ocean dimethyl sulfide climatology estimated from observations and an
artificial neural network.}
\textit{Biogeosciences}, 17, 5335--5354, 2020.

\end{thebibliography}

\end{document}

